\chapter{The taxonomy and where it comes from}
\label{ch:taxonomy}

\section{Provenance}

The ground truth label space is the WHATWG autofill field name token set
\cite{whatwg-autofill} plus a small, enumerated set of extra labels. Nothing in
it was invented for convenience, and the extras are named individually rather
than admitted by a rule, because a category that can be added by a rule is a
category that grows until the label space is nobody's specification.

The extras exist for outcomes the token set has no word for, and each one earns
its place by naming something a report has to be able to say:

\begin{itemize}
  \item \code{NOT\_AUTOFILLABLE}, for a control that is a real control and has
    no autofill meaning at all: a search box, a consent checkbox, a quantity
    stepper. Without it the only honest answer for such a field would be
    \code{UNKNOWN}, which reads as \emph{the tool could not tell} when the truth
    is \emph{there is nothing to tell}.
  \item \code{UNKNOWN}, for a control whose evidence genuinely does not decide.
    This is an answer, not an absence of one.
  \item \code{COMPOSITE\_UNSPLIT}, for one input holding what the specification
    models as several fields, the archetype being a single expiry box
    accepting month and year together.
  \item \code{CC\_EXP\_SPLIT\_MONTH} and \code{CC\_EXP\_SPLIT\_YEAR}, for the
    two halves of a split expiry pair, which are individually meaningless and
    jointly a \code{cc-exp}.
\end{itemize}

\section{The growth rule, made mechanical}

A label may exist only if it is a specification token or an enumerated extra,
\emph{and} is emitted by at least one corpus answer key, \emph{and} is reachable
by at least one rule or present in the model's training distribution, \emph{and}
is asserted by at least one test. A label that cannot satisfy all four is
deleted, or carries a dated and tracked exemption.

That is a rule somebody will forget, so it is a script.
\code{scripts/check\_reachability.py} runs in continuous integration, reads the
taxonomy from the one source file that defines it, and fails the build on any
clause it can currently enforce. Its clauses were activated one at a time as the
phases that made each enforceable landed, and each activation moved a clause out
of the file's own pending list in the same commit that made it checkable, which
is the only way a pending list stays honest.

Two further properties of the check are worth naming because they are
strengthenings rather than restatements.

\textbf{It scans for label strings written outside the taxonomy module.} Ground
rule 6 says the taxonomy has exactly one definition in code. The check makes
that mechanical by walking the source and the scripts for label literals. It
deliberately excludes the handful of tokens that are also ordinary English
words, because a check that fires on the word \code{name} in a comment is a
check that gets switched off within a week.

\textbf{It requires a training floor.} Every label a model can predict must
appear in the training partition with at least a full template's worth of rows.
That clause was added at the corpus repair described in
Chapter~\ref{ch:corpus}, after three labels were found sitting at zero training
rows, which means the model had been asked about classes it had never once seen
and could only ever get wrong.

\section{Coverage in the corpus as generated}

Every label in the taxonomy is emitted by at least one answer key in the
generated corpus. Table~\ref{tab:label-coverage} in Appendix~\ref{app:tables}
lists the keyed field count per label, which is the evidence for clause (b) of
the growth rule and also the honest picture of how uneven the class distribution
is. Several labels are carried by a single template in a single position, and
the per label results in Table~\ref{tab:per-label} should be read against that
column rather than on their own.

\section{Why not a bespoke ontology}

The alternative was tempting and is worth refusing in writing.

A bespoke ontology would have let the corpus, the rules and the model agree with
each other perfectly, because all three would have been designed together. It
would also have made the tool's central output, the exact attribute value to
add, a translation step rather than a formatting step, and every disagreement
between the ontology and the specification would have become a silent wrong
answer in production markup.

More importantly, it would have made the evaluation meaningless outside this
repository. Measuring a classifier against categories the same project invented
measures internal consistency. Measuring it against the token set that browsers
actually implement measures the thing the tool claims to do.
